% Shared local evidence, measured 10 September 2026.
\clearpage
\cacheheading{FORM cache measurements with character projection prefilled}
The controlled comparison used nine theories at $t^{12}$ and $t^{18}$, one
warmup and five measured repetitions per mode, with fresh clients and one worker.
Both decompositions and final singlet coefficients were already filled.
The character database was query-only and guards prohibited LiE execution.
All 270 timed indices agreed exactly. Medians below are full calculation times
in seconds, excluding setup and Python/Sage startup.

\begin{tabularx}{\textwidth}{@{}Yrrrr@{}}
\toprule Theory, through $t^{18}$ & Disabled & Miss/save & Hit & Speedup\\\midrule
$SU(2)$, 4 fundamentals & 0.4323 & 0.4916 & 0.0617 & $7.00\times$\\
$SU(3)$, 6 fundamentals & 1.7220 & 1.7959 & 0.2304 & $7.48\times$\\
$SU(5)$, 10 fundamentals & 1.7265 & 1.7839 & 0.2648 & $6.52\times$\\
$SU(3)$, 1 adjoint ($\mathcal N=4$) & 0.2808 & 0.3180 & 0.0429 & $6.54\times$\\
$Sp(2)$, 6 fundamentals & 0.4434 & 0.4817 & 0.0634 & $7.00\times$\\
$\operatorname{Spin}(8)$, 6 vectors & 0.4329 & 0.4848 & 0.0655 & $6.60\times$\\
$G_2$, 4 fundamentals & 0.4418 & 0.4874 & 0.0640 & $6.90\times$\\
$SU(2)^2$, 2 bifundamentals & 2.7421 & 2.7843 & 0.3251 & $8.44\times$\\
$SU(5)$, symmetric + antisymmetric & 13.7565 & 14.1896 & 1.8771 & $7.33\times$\\\bottomrule
\end{tabularx}
All matter in this table is full hypermultiplets; $Sp(2)=USp(4)$. Disabled mode
bypasses FORM SQLite and executes/parses FORM; miss/save begins with an
initialized empty FORM database and includes serialization and insertion.
Hit mode decodes an existing expansion with FORM execution forbidden. The
speedup is disabled time divided by hit time. At $t^{12}$ the range was
$3.23$--$6.92\times$. These are warm-filesystem measurements on this machine,
not cold-disk or contention benchmarks. Evidence: \code{report.md},
\code{results.json} and \code{benchmark.py} under
\code{output/benchmarks/form_expansion_theories/}.

\textbf{Cross-theory hits.} Two pairs share byte-identical programs:
$SU(2)$ with four fundamentals and $G_2$ with four fundamentals;
$Sp(2)$ with six fundamentals and $\operatorname{Spin}(8)$ with six vectors.
Both directions at both orders gave eight passing reuse checks. The first
theory executed FORM once; a new client for the second executed/parsed FORM
zero times, decoded once, and left the database at one row. Each final index
matched its own independent fresh-FORM result, while the paired final indices
differed. A changed-program control missed as expected. See
\code{cross_theory_reuse.md}, \code{cross_theory_reuse.json} and
\code{check_cross_theory_reuse.py} in that same directory.

\clearpage
\cacheheading{I/O cost, decomposition planning and builder verification}
A separate I/O benchmark used the 113,831-term $t^{18}$ expansion: a 589-byte
program and 5,907,568-byte JSON payload (5.91 MB). Thirty measured repetitions
followed three warmups, using SQLite WAL with synchronous FULL on the project
filesystem. Medians are milliseconds:
\begin{center}\begin{tabular}{lr}
\toprule Operation & Median (ms)\\\midrule
Serialize exact terms & 179.32\\
SQLite insert and commit & 42.07\\
Total save, excluding database initialization & 221.51\\
SQLite read and fetch JSON & 3.08\\
Reconstruct \code{IndexFormTerm} objects & 866.42\\
Hit with an already open connection & 870.81\\
Hit with a new client, including close & 894.65\\\bottomrule
\end{tabular}\end{center}
Initialization separately took 31.43 ms. No FORM execution/parsing is included;
every round trip was checked exactly. Python reconstruction dominated loading.
Individual component medians need not add to the median total. This is a
single-process warm-filesystem test, not concurrent lock contention.
Evidence: \code{output/benchmarks/form_expansion_cache_io/report.md},
\code{timings.json} and \code{benchmark.py}.

\textbf{Cached-factor planner.} The temporary candidate matched 130 products
over $A_1$, $A_2$ fundamental/adjoint, $C_2$ and $G_2$ in serial and three-worker
runs; $SU(2)$ cases also matched an independent weight oracle. Prepared-cache
misses improved $1.31$--$1.72\times$ for the single-request API and
$1.03$--$1.17\times$ for the batch API, reducing two tensor calls to one. A case
with only composite entries also reduced one Adams call to zero. However, 72
full-index comparisons over six theories, with FORM prefilled and either a cold
character cache or a cache filled through $t^{12}$, showed \emph{no consistent
overall speedup}. All indices agreed and tensor-call counts were unchanged.
The retained baseline, candidate, report, scripts and raw measurements are in
\code{output/benchmarks/adams_cache_planner/}.

\textbf{Complete builder.} Twelve regressions cover all 29 $SU(2)$ products
through order six against independent weights, serial/parallel agreement for
$A_1,A_2,C_2,G_2$ through order five, pool reuse, warm skips, partial-cache resume,
failure persistence, input validation and both CLI entry points. A real worker
log records exactly one Adams call at each of orders two through six and 23
tensor calls, with multiple worker PIDs. This verifies the operation-count claim
on the tested build; no broad prebuild speedup has been measured.

The full suite rerun during this documentation update ran \textbf{209 tests: 207 passed and
two live-MySQL tests skipped}. It uses actual Sage/FORM/LiE for backend tests;
MySQL migration and insertion unit tests use mocks unless a test server is
explicitly configured. The new modules are
\code{test_form_expansion_cache.py}, \code{test_cached_decomposition_planning.py}
and \code{test_build_decomposition_cache.py}. These measurements and algorithms
are source-derived project evidence, not results attributed to external papers.
